\subsection{A positive density of nonsingular agreements gives linear MCA}
\label{sec:nonsingular-agreement-mca}

The preceding component counts do not by themselves bound all bad
labels linearly. Conditioning Taylor reconstruction at an agreeing
coordinate gives a different sufficient condition that also covers
surface components. It isolates singular agreement coordinates as
the remaining obstruction to this argument.

\begin{theorem}[Nonsingular-agreement bound]
\label{thm:nonsingular-agreement-mca}
Let $F$ have characteristic zero or $p>D$, with $D\ge1$, and let
$Q(X,z,u,v)\ne0$ have jet degree at most $B\ge1$ and challenge
degree at most $H$. Its $X$-degree is unrestricted. Put
\[
 \tau=\max\{0,2D-3\},\quad b=1+\tau(B-1),\quad
 T_0=b+\tau H,\quad K=H+BT_0,\quad q=B+H.
\]
Fix a received line $f+zg$ on $n$ distinct coordinates, an agreement
threshold $D<A\le n$, and an integer $L\ge1$.
Consider labels admitting a degree-$\le D$ polynomial solution
$Q(X,z,P,P')=0$ whose full agreement support $S$ satisfies
$|S|\ge A$, has no degree-$\le D$ polynomial agreeing with $g$ on
all of $S$, and contains at least $L$ coordinates $x$ with
\[
 Q_v(x,z,P(x),P'(x))\ne0.
\]
The number of these labels is at most
\begin{equation}
 \frac nL\left(qK+qT_0\frac{n-D}{A-D}+qn\right).
 \label{eq:nonsingular-agreement-bound}
\end{equation}
In particular it is $O_{B,H,\eta,\lambda}(n)$ when
$A-D\ge\eta n$ and $L\ge\lambda n$, with $\eta,\lambda>0$ fixed.
\end{theorem}
\begin{proof}
Fix a domain point $x$ and substitute $u=f_x+zg_x$ into the initial
equation. The resulting polynomial in $(z,v)$ has degree at most $q$.
If it vanishes identically, so does its $v$ derivative, and there
are no nonsingular initial data to count at this coordinate.
Otherwise work on the open locus
$S_0=Q_v(x,z,f_x+zg_x,v)\ne0$ of this plane curve.

The Taylor reconstruction in
Theorem~\ref{thm:actual-first-order-components} applies at the fixed
center $x$ on this open locus. Its common denominator is $S_0^\tau$,
of degree at most $T_0-1$; its coefficient numerators and $zS_0^\tau$
have degree at most $T_0$. Thus its projective coordinate degree is
at most $T_0$. Exact solutions are selected by residual equations
of degree at most $K$. Reconstruction and the initial-jet map are
inverse regular maps on this open exact-solution locus.

The isolated points of the residual locus number at most $qK$.
Indeed, on every component of the initial plane curve not wholly
satisfying the residuals, some residual is nonzero. A generic linear
combination over an algebraic closure is nonzero on each such component;
B\'ezout gives the bound. Components contained in $S_0=0$ are omitted,
and points on a retained curve are accounted for with that curve.

The retained source curves are components of the initial plane
equation, so there are at most $q$ and their degrees sum to at most
$q$. Their reconstructed image curves have total degree at most
$qT_0$, by pulling back a generic hyperplane under the rational map
of degree at most $T_0$. Initial-jet recovery ensures none of these
images collapses to a point.

For every nonaffine image curve, the incidence argument of
Corollary~\ref{cor:nonaffine-actual-curves} gives at most its degree
times $(n-D)/(A-D)$ nearby pairs. That argument applies to these
curves whether or not they are components of the full joint locus.
The total contribution is at most $qT_0(n-D)/(A-D)$.

Each remaining image curve is a graph $P=F_0+zG_0$ and has at most
$n$ bad labels. A support consisting only of persistent agreements
has the common witnesses $F_0,G_0$; a bad support therefore contains
an accidental agreement at a nonpersistent coordinate, which forces
one label. If the graph does not descend to $F$, it has at most one
$F$-rational point at distinct challenge labels: two such points
would recover $F_0,G_0$ over $F$. This also satisfies the bound $n$.
There are at most $q$ graphs.

Hence this fixed coordinate occurs nonsingularly in a bad nearby
witness for at most $qK+qT_0(n-D)/(A-D)+qn$ labels. Select one
qualifying witness per label. Each supplies at least $L$ such
coordinate--label incidences. Summing over all $n$ coordinates and
dividing by $L$ proves~\eqref{eq:nonsingular-agreement-bound}.
\end{proof}

\begin{corollary}[Value-independent separant]
\label{cor:value-independent-separant}
Suppose
\[
 Q(X,z,u,v)=R(X,z)v-\mathcal A(X,z,u),\qquad R\ne0,
\]
under the same degree and characteristic hypotheses, and put
$r=\deg_XR$. If $A>\max\{D,r\}$, all its full-support bad labels
number at most
\begin{equation}
 H+\frac n{A-r}\left(qK+qT_0\frac{n-D}{A-D}+qn\right).
 \label{eq:quasilinear-linear-mca}
\end{equation}
This is $O(n)$ for fixed $B,H$ and fixed positive lower bounds on
$(A-D)/n$ and $(A-r)/n$. The value degree of $\mathcal A$ is arbitrary.
At each fixed label with $R(X,z)\not\equiv0$, its ordinary list size
at threshold $A$ is at most $\lfloor(n-r)/(A-r)\rfloor$.
\end{corollary}
\begin{proof}
There are at most $H$ labels with $R(X,z)\equiv0$, using any one
nonzero $X$-coefficient of $R$. At all others, at most $r$ domain
points are zeros of $R$, so every nearby candidate has at least
$A-r$ nonsingular agreements. Apply the theorem with $L=A-r$.

For the ordinary list bound, two distinct solutions at the same
label have a nonzero difference $\Delta$ satisfying
$R\Delta'=C(X)\Delta$. At a zero of $\Delta$ where $R$ is nonzero,
the multiplicity $1\le m\le D<p$ makes the left side have order
$m-1$ and the right side order at least $m$, a contradiction.
Thus a value at a nonsingular coordinate determines at most one
solution. If $r_z\le r$ domain points are singular, nearby candidates
have disjoint regular agreement supports of size at least $A-r_z$.
Their number is at most $(n-r_z)/(A-r_z)\le(n-r)/(A-r)$.
\end{proof}

\paragraph{Sharp linear order in a nonlinear subclass.}
For a nonzero polynomial $W$, the equation
\[
 WP'-W'P-P(P-W)=0
\]
has exactly the polynomial solutions $0,W$. To see this, write
$Z=P/W$, giving $Z'=Z(Z-1)$. If $Z\ne0$, then $U=1-1/Z$ satisfies
$U'=U$. In characteristic zero, a nonzero rational logarithmic
derivative $U'/U$ vanishes at infinity and cannot equal one. In
characteristic $p$, the $p$-fold ordinary derivative is zero on
$F(X)$, whereas $U'=U$ would make it equal $U$. Hence $U=0$.

Let $n=12s$, $D=3s$, $A=6s$, and choose a prime $p>12s$. On
domain $\{0,\ldots,n-1\}$ let $W$ be the locator of the first $3s$
points. Put $f=g=0$ on these roots and on a second block of $3s$
points; put $f=W,g=0$ on a third block of $3s$. At the final $3s$
coordinates, indexed by $j=1,\ldots,3s$, put $f=3jW,g=W$.
Each candidate $cW$, $c\in\{0,1\}$, has $6s$ persistent agreements
and acquires one more at each label $z=c-3j$. These $6s=n/2$ labels
are distinct modulo $p$. On each resulting support, $g$ has $6s>D$
zeros and one nonzero value, so it has no degree-$\le D$ explanation.
Every other label has only persistent supports and is good. Thus
exactly $n/2$ labels are bad, with jet degree two, challenge degree
zero, code rate tending to $1/4$, and agreement gap $1/4$.

\paragraph{Scope and checks.}
The theorem does not assert that interpolation equations have many
nonsingular agreement coordinates. At fixed gap its bound is
$O(n^2/L)$; a superlinear example must concentrate on witnesses with
a vanishing fraction of such coordinates. The general first-order
conjecture remains open. The checks in
\path{research/quasilinear_first_order/} enumerate small-field
polynomial solutions, verify regular-coordinate uniqueness and list
incidences, include implicit equations with actual solution surfaces,
and check the exact $n/2$ lower construction at five lengths. These
finite checks supplement the reconstruction and incidence proof;
no independent review or literature-priority claim is made.
